\section{Modely a trénování}

\subsection*{Baseline}

Jako referenční modely slouží čtyři přístupy pracující jen s metadaty. Nejjednodušší baseline
předpovídá pro každou hru průměr $y$ z trénovací množiny. Dále je použita lineární regrese
a gradient boosting (\texttt{HistGradientBoostingRegressor} s výchozím nastavením) z knihovny
scikit-learn~\cite{pedregosa2011sklearn}. Poslední baseline je neuronová síť MLP se skrytou
vrstvou o 128 neuronech a aktivací ReLU (10\,113 parametrů). Ta je zároveň referencí pro otázku,
zda obrázky přinášejí informaci navíc.

\subsection*{Hybridní architektura}

Modely A a B sdílejí stejnou hybridní architekturu (obrázek~\ref{fig:architektura}) a liší se
pouze encoderem obrázků. Encoder $e_\psi$ převede každý ze čtyř obrázků na vektor délky $d$.
Vektory se zprůměrují, spojí s metadaty a MLP hlava $g_\phi$ z nich spočítá predikci:
\begin{equation}
    \hat{y}_i = g_\phi\left(\left[\frac{1}{4} \sum_{k=1}^{4} e_\psi(I_{i,k}) \,;\, \mathbf{m}_i\right]\right),
\end{equation}
kde $[\,\cdot\,;\,\cdot\,]$ značí spojení vektorů. Jeden encoder zpracovává všechny obrázky,
takže se počet jeho parametrů s počtem obrázků nezvyšuje. Průměrování navíc nezávisí na pořadí
obrázků. Hlava $g_\phi$ má stejnou strukturu jako baseline MLP: lineární vrstvu s $h$ výstupy,
ReLU a lineární výstupní vrstvu s jedním neuronem. Šířka $h$ je hyperparametr (kapitola~4),
ve finálních konfiguracích obou modelů je $h = 64$.

\begin{figure}[h]
    \centering
    \begin{tikzpicture}[
        box/.style={draw, rounded corners, align=center, minimum height=8mm, inner sep=3pt, font=\small},
        arr/.style={-Stealth}
    ]
        % čtyři obrázky hry a sdílený encoder
        \foreach \i/\name in {1/header, 2/screenshot 1, 3/screenshot 2, 4/screenshot 3} {
            \node[box, minimum width=2.3cm] (img\i) at (0, {-(\i-1)*1.1}) {\name};
            \node[box, fill=blue!10, minimum width=1.8cm] (enc\i) at (2.9, {-(\i-1)*1.1}) {$e_\psi$};
            \draw[arr] (img\i) -- (enc\i);
        }
        \node[draw, dashed, rounded corners, inner sep=5pt, fit=(enc1)(enc4),
              label={[font=\small]above:encoder (sdílené váhy)}] {};

        % průměr vektorů obrázků a spojení s metadaty
        \node[box] (mean) at (5.6, -1.65) {průměr\\$\mathbb{R}^{d}$};
        \foreach \i in {1,...,4} \draw[arr] (enc\i.east) -- (mean.west);
        \node[box] (meta) at (5.6, -4.0) {metadata\\$\mathbb{R}^{77}$};
        \node[box] (cat) at (8.0, -2.8) {spojení\\$\mathbb{R}^{d+77}$};
        \draw[arr] (mean) -- (cat);
        \draw[arr] (meta) -- (cat);

        % MLP hlava a výstup
        \node[box, fill=orange!15] (head) at (10.9, -2.8) {hlava $g_\phi$\\Linear($h$), ReLU\\Linear(1)};
        \draw[arr] (cat) -- (head);
        \node (out) at (13.0, -2.8) {$\hat{y}$};
        \draw[arr] (head) -- (out);
    \end{tikzpicture}
    \caption{Hybridní architektura modelů A a B. Sdílený encoder převede každý obrázek na vektor
    délky $d$ ($d = 64$ pro model A, $d = 768$ pro model B), vektory se zprůměrují a spojí s metadaty.}
    \label{fig:architektura}
\end{figure}

\subsection*{Model A: vlastní CNN}

Encoder modelu A je konvoluční síť trénovaná od náhodné inicializace. Skládá se ze tří bloků,
každý obsahuje dvě konvoluce $3 \times 3$ s aktivací ReLU a max pooling $2 \times 2$, který
zmenšuje prostorové rozlišení na polovinu. Počet kanálů se v blocích zdvojnásobuje
(16, 32, 64), rozlišení klesá ze $128 \times 128$ na $16 \times 16$. První bloky zachycují
jednoduché vzory (hrany, barvy, textury), další bloky z nich skládají složitější struktury.
Global average pooling zprůměruje každou ze 64 map příznaků na jedno číslo. Výsledný vektor
tak popisuje, \emph{zda} se vzor v obrázku vyskytuje, ale ne \emph{kde}. Následují dvě plně
propojené vrstvy (64 $\rightarrow$ 128 $\rightarrow$ 64) s aktivací ReLU, takže $d = 64$.

Global average pooling nahradil v původní verzi sítě natažení map příznaků do plně propojené
vrstvy s přibližně 4,2 milionu parametrů, se kterou se síť rychle přeučovala.

\subsection*{Model B: předtrénovaný encoder}

Model B používá jako encoder síť ConvNeXt-Tiny~\cite{liu2022convnext} předtrénovanou
na datasetu ImageNet~\cite{deng2009imagenet} z knihovny torchvision. Klasifikační vrstva
pro 1000 tříd je odstraněna, výstupem je vektor délky $d = 768$. Váhy encoderu jsou zmrazené,
trénuje se pouze hlava $g_\phi$ (linear probing s MLP hlavou). Obrázky se normalizují
statistikami ImageNetu. Encoder zůstává i během trénování v režimu inference, aby se nepoužíval
stochastic depth, který náhodně vynechává bloky sítě. Síť byla předtrénována na obrázcích
$224 \times 224$, zde dostává obrázky $128 \times 128$.

\begin{table}[h]
    \centering
    \begin{tabular}{lrrr}
        \toprule
        Model & Encoder & Hlava & Trénované \\
        \midrule
        MLP (metadata) & -- & 10\,113 & 10\,113 \\
        Model A & 88\,656 & 9\,153 & 97\,809 \\
        Model B & 27\,820\,128 & 54\,209 & 54\,209 \\
        \bottomrule
    \end{tabular}
    \caption{Počty parametrů modelů ve finálních konfiguracích (hlava modelů A a B s $h = 64$).}
    \label{tab:parametry}
\end{table}

\subsection*{Trénování}

Všechny neuronové sítě jsou implementovány v PyTorch~\cite{paszke2019pytorch} a trénovány
stejným postupem. Minimalizuje se Huberova ztráta (kapitola~1) optimizerem
AdamW~\cite{loshchilov2019adamw}, který aktualizuje parametry podle
\begin{equation}
    \theta_t = \theta_{t-1} - \eta \left( \frac{\hat{m}_t}{\sqrt{\hat{v}_t} + \epsilon} + \lambda \theta_{t-1} \right),
\end{equation}
kde $\hat{m}_t$ a $\hat{v}_t$ jsou korigované klouzavé průměry gradientu a jeho druhé mocniny,
$\eta$ je learning rate a $\lambda$ weight decay. Na rozdíl od L2 regularizace v původním
Adamu se weight decay neškáluje adaptivním krokem. MLP nad metadaty používá $\eta = 10^{-3}$,
$\lambda = 10^{-4}$ a velikost dávky 256. Hyperparametry modelů A a B jsou vybrány random
searchem (kapitola~4).

Po každé epoše se spočítá MAE na validační množině a ukládá se model z nejlepší epochy.
Trénování trvá nejvýše 15 epoch a končí dříve (early stopping), pokud se validační MAE
nezlepší 3 epochy po sobě. Dopředný průchod probíhá ve smíšené přesnosti (bfloat16), ztráta
se počítá ve float32. Formát bfloat16 má stejný rozsah exponentu jako float32, takže není
potřeba škálovat gradienty. Každý model je natrénován se třemi seedy
(123, 420, 666), které určují inicializaci vah i pořadí trénovacích dat.